%!TEX root =  tfg.tex
\chapter{Costes}

\begin{quotation}[Escritores y Raperos]{Jim Johnston y Naughty by Nature}
Here Comes the Money\end{quotation}

\begin{abstract}
En este capítulo se realiza el desglose de los costes asociados al proyecto, basándose en los salarios reales de puestos tecnológicos.
\end{abstract}

\section{Resumen de costes del proyecto}

\begin{table*}[h!]
	\centering
	\begin{coolTable}{ll}{2}
{Resumen del proyecto}
	\textbf{Costes de personal}&14.601,17 \euro\\
	Desarrollador Full-Stack&9.522,50 \euro\\
	Becario Documentaci\'{o}n&5.078,67 \euro\\
	\textbf{Costes materiales}&67,23 \euro\\
	\textbf{Costes indirectos}&1.466,84 \euro\\
	\midrule
	\textbf{TOTAL}&16.135,24 \euro\\
	\end{coolTable}
	\caption{Tabla resumen de costes del proyecto}
\end{table*}

\section{Costes de personal}

Para el cálculo de los costes de personal se han consultado los datos de la \textbf{Guía Salarial 2026} de Manfred\footnote{\url{https://www.getmanfred.com/blog/guia-salarial-2026-salarios-en-tecnologia-espana-manfred}}, que recopila información salarial de más de 120.000 profesionales de tecnología en España.

Dado que el desarrollo de Scubex ha requerido asumir múltiples roles y tareas, se han considerado dos perfiles:

\subsection{Desarrollador Full-Stack (Arquitectura e Implementación)}

Se ha considerado como referencia el salario de \textbf{Desarrollador Full-Stack} con experiencia de 5-10 años, cuyo rango salarial se sitúa entre 41.000-50.000\euro{} anuales según dicha guía. Se ha tomado un valor medio de \textbf{45.000\euro{} brutos anuales}.

\subsubsection{Cálculo proporcional al tiempo dedicado}

La carga de trabajo estimada para el TFG es de \textbf{293 horas}, frente a una jornada laboral anual estándar de \textbf{1.800 horas}. Por tanto, el coste proporcional se calcula como:

\begin{equation}
	\text{Salario proporcional} = 45.000\text{\euro} \times \frac{293}{1.800} = 7.325\text{\euro}
\end{equation}

\subsubsection{Costes sociales}

A este salario bruto se le añaden los \textbf{costes sociales} (Seguridad Social, contingencias comunes, desempleo, formación profesional, etc.) que suponen aproximadamente el \textbf{30\%} del salario bruto:

\begin{equation}
	\text{Costes sociales} = 7.325\text{\euro} \times 0.30 = 2.197,50\text{\euro}
\end{equation}

\subsubsection{Total Desarrollador Full-Stack}

\begin{equation}
	\text{Total} = 7.325\text{\euro} + 2.197,50\text{\euro} = 9.522,50\text{\euro}
\end{equation}

\subsection{Becario Junior (Redacción de Documentación)}

Para la redacción de la memoria académica y documentación técnica se ha considerado el perfil de un \textbf{becario junior} con experiencia menor a 2 años. Según la guía salarial, los desarrolladores junior se sitúan en el rango de 20.000-30.000\euro{} anuales. Se ha tomado un valor de \textbf{24.000\euro{} brutos anuales}.

\subsubsection{Cálculo proporcional al tiempo dedicado}
De la misma forma que se ha calculado anteriormente:

\begin{equation}
	\text{Salario proporcional} = 24.000\text{\euro} \times \frac{293}{1.800} = 3.906,67\text{\euro}
\end{equation}

\subsubsection{Costes sociales}

\begin{equation}
	\text{Costes sociales} = 3.906,67\text{\euro} \times 0.30 = 1.172\text{\euro}
\end{equation}

\subsubsection{Total Becario}

\begin{equation}
	\text{Total} = 3.906,67\text{\euro} + 1.172\text{\euro} = 5.078,67\text{\euro}
\end{equation}

\subsection{Total costes de personal}

\begin{equation}
	\text{Total} = 9.522,50\text{\euro} + 5.078,67\text{\euro} = 14.601,17\text{\euro}
\end{equation}

\section{Costes materiales}

\subsection{Amortización de equipamiento informático}

Según el criterio de amortización fiscal establecido por la Agencia Tributaria española, los equipos informáticos se amortizan a un \textbf{máximo del 25\% anual} (vida útil de 4 años).

Para el desarrollo del proyecto se ha utilizado un portátil con un coste de adquisición de \textbf{1.200\euro}:

\begin{itemize}
	\item \textbf{Amortización anual}: $1.200\text{\euro} \times 0.25 = 300\text{\euro}$
	\item \textbf{Amortización proporcional} (293h / 1.800h): $300\text{\euro} \times \frac{293}{1.800} = 48{,}83\text{\euro}$
\end{itemize}

\subsection{Despliegue en la nube}

El proyecto se encuentra desplegado en producción y preproducción mediante dos servicios:

\begin{itemize}
	\item \textbf{Railway} (backend Spring Boot + PostgreSQL): Plan Hobby a 5\,USD/mes (\textasciitilde{}4,60\,\euro/mes). Este plan incluye 5\,USD de créditos de uso mensuales; superado ese umbral, se factura por consumo real de recursos. Actualmente debido al tráfico bajo--moderado, el crédito mensual de 5\,USD resulta suficiente.
	\item \textbf{Vercel} (frontend React/Vite): Plan Hobby gratuito, sin coste.
\end{itemize}

\textbf{Coste estimado de despliegue durante el desarrollo del TFG} (4 meses):

\begin{equation}
	\text{Railway} = 4,60\text{\,\euro/mes} \times 4\text{\,meses} = 18,40\text{\,\euro}
\end{equation}

\begin{equation}
	\text{Total despliegue} = 18,40\text{\,\euro} + 0\text{\,\euro\,(Vercel)} = 18,40\text{\,\euro}
\end{equation}

En un escenario de producción real con mayor tráfico, los costes de Railway podrían aumentar proporcionalmente al uso de CPU, RAM y almacenamiento consumidos más allá del crédito incluido. En este caso podría crearse un plan personalizado, que incluya las necesidades concretas de la aplicación.

\subsection{Total costes materiales}

\begin{equation}
	\text{Total} = 48{,}83\text{\,\euro} + 18{,}40\text{\,\euro} = 67{,}23\text{\,\euro}
\end{equation}

\section{Costes indirectos}

Los costes indirectos incluyen gastos generales asociados al desarrollo del proyecto que no se imputan directamente: electricidad, conexión a internet, espacio de trabajo, cuotas mensuales de asistentes de IA, etc.

Se ha tomado un valor medio del \textbf{10\%} de la suma de costes de personal y materiales:

\begin{equation}
	\text{Costes indirectos} = (14.601,17\text{\,\euro} + 67,23\text{\,\euro}) \times 0.10 = 1.466,84\text{\,\euro}
\end{equation}

\section{Coste total del proyecto}

\begin{table*}[h!]
	\centering
	\begin{coolTable}{lr}{2}
{Desglose de costes}
	\textbf{Concepto}&\textbf{Importe}\\
	\midrule
	Costes de personal&14.601,17\euro\\
	\quad Desarrollador Full-Stack&9.522,50\euro\\
	\quad Becario Documentaci\'{o}n&5.078,67\euro\\
	Costes materiales&67,23\euro\\
	\quad Amortización equipamiento&48,83\euro\\
	\quad Despliegue Railway (4 meses)&18,40\euro\\
	Costes indirectos&1.466,84\euro\\
	\midrule
	\textbf{TOTAL PROYECTO}&\textbf{16.135,24\euro}\\
	\end{coolTable}
	\caption{Coste total del proyecto Scubex}
\end{table*}